Todo lo que pasa en un espejo plano se explica con una sola regla, la ley de la reflexión. Si trazamos una línea perpendicular al espejo en el punto donde llega un rayo, a la que llamamos normal, el ángulo entre el rayo que llega y la normal es igual al ángulo entre el rayo que sale y la normal:
\begin{equation}
\theta_i = \theta_r .
\label{eq:reflexion}
\end{equation}
Esta ley, como vimos en la práctica del principio de Fermat, es la consecuencia de que la luz siga el camino más corto.

\subsection{Dónde se forma la imagen}
Pensemos en un punto luminoso O colocado a una distancia $s_o$ frente al espejo. De O salen rayos en todas direcciones. Uno de ellos llega perpendicular al espejo y rebota sobre sí mismo. Otro llega con un ángulo $\theta$ y sale con ese mismo ángulo. Si prolongamos hacia atrás los dos rayos reflejados, se cruzan en un punto I detrás del espejo, a una distancia $s_i$.

Ahora viene la parte bonita. El rayo inclinado toca el espejo a una altura $h$ por encima del rayo perpendicular, y con la geometría de los dos triángulos rectángulos que se forman se tiene
\begin{equation}
\tan\theta = \frac{h}{s_o} \qquad \text{y} \qquad \tan\theta = \frac{h}{s_i} ,
\label{eq:tangentes}
\end{equation}
porque el ángulo es el mismo de los dos lados gracias a la ecuación \eqref{eq:reflexion}. Al igualar ambas expresiones queda
\begin{equation}
s_i = s_o .
\label{eq:distancias}
\end{equation}
Este resultado no depende de $\theta$, así que todos los rayos reflejados parecen salir del mismo punto I. Ese punto es la imagen.

\subsection{Características de la imagen}
La imagen de un espejo plano tiene cuatro rasgos que vale la pena tener presentes:
\begin{itemize}
\item Es virtual, porque los rayos no pasan realmente por ella, solo sus prolongaciones.
\item Está a la misma distancia detrás del espejo que el objeto delante, como dice la ecuación \eqref{eq:distancias}.
\item Tiene el mismo tamaño que el objeto.
\item Está invertida de izquierda a derecha, o mejor dicho, de adelante hacia atrás.
\end{itemize}
El tercer punto se resume con el aumento lateral, que compara la altura de la imagen $y_i$ con la del objeto $y_o$:
\begin{equation}
M = \frac{y_i}{y_o} = \frac{s_i}{s_o} = 1 .
\label{eq:aumento}
\end{equation}

\subsection{Método de paralaje}
Como la imagen es virtual, no se puede atrapar en una pantalla, y hay que encontrarla de otra manera. El truco es el paralaje, el mismo efecto que hace que un objeto cercano parezca moverse respecto a uno lejano cuando movemos la cabeza. Si colocamos un segundo objeto detrás del espejo y movemos la cabeza de un lado a otro, el objeto y la imagen se desplazan uno respecto al otro mientras estén a distancias distintas de nuestros ojos, y dejan de hacerlo cuando están en el mismo lugar. Ese es el punto donde está la imagen.
